% app-and-vc-lifecycle.tex — UIViewController callbacks in exact order + app/scene states.
% Sources (checked 2026-10-05):
%   developer.apple.com/documentation/uikit/uiviewcontroller (loadView: do not call super when
%     overriding; viewIsAppearing(_:) — iOS 13+ with the iOS 17 SDK; containment: addChild,
%     didMove(toParent:), willMove(toParent:), removeFromParent;
%     shouldAutomaticallyForwardAppearanceMethods default true)
%   developer.apple.com/documentation/uikit/managing-your-app-s-life-cycle (scene-based states;
%     app-delegate state callbacks not called once scenes are adopted)
%   developer.apple.com/documentation/uikit/uiscenedelegate · .../swiftui/scenephase
%   developer.apple.com/documentation/xcode/addressing-watchdog-terminations (0x8badf00d)
% Not documented by Apple as a number: the ~20 s launch watchdog (marked \unverified).
% @labels: area=ios-swift kind=concept platform=ios topic=ui,platform-apis discussion=58
% @tags: uiviewcontroller, loadview, viewdidload, viewisappearing, viewwillappear, viewdidlayoutsubviews, app-states, scene-delegate, uiapplicationdelegate, scenephase, cold-launch, jetsam
\documentclass[8pt]{extarticle}
\usepackage{printup-sheet}
\usepackage{array}

\lstdefinelanguage{SwiftSheet}{
  morekeywords={protocol,class,final,struct,enum,func,var,let,weak,init,deinit,
    override,super,if,else,return,guard,self,nil,try,await,async,throws,private,some,
    true,false,AnyObject,Void,String,Bool},
  sensitive=true, morecomment=[l]{//}, morecomment=[s]{/*}{*/}, morestring=[b]"}

\tikzset{
  lc/.style={box, font=\scriptsize, inner sep=1pt, minimum height=8mm, text width=12.6mm},
  once/.style={lc, draw=sheetBlue, fill=sheetBlue!10},
  app/.style={lc, draw=sheetGreen!70!black, fill=sheetGreen!12},
  lay/.style={lc, draw=sheetOrange, fill=sheetOrange!12},
  dis/.style={lc, draw=sheetRed, fill=sheetRed!8},
  lbl/.style={font=\tiny, text=black!75, inner sep=1pt, align=center},
  st/.style={box, font=\scriptsize\bfseries, minimum width=15mm, minimum height=6mm, inner sep=1.5pt},
  ed/.style={font=\tiny\ttfamily, text=black!80, inner sep=1pt, align=center},
  br/.style={decorate, decoration={brace, amplitude=3pt}, thick},
}

\begin{document}

\sheettitle{App \& view-controller lifecycle — exact order}{ios-swift · folio}

\oneliner{A VC's view is \textbf{lazy}: first access to \texttt{.view} runs
\texttt{loadView} then \texttt{viewDidLoad} (\textbf{once}); every show runs
\texttt{viewWillAppear} → \texttt{viewIsAppearing} → layout → \texttt{viewDidAppear}; every
hide runs \texttt{viewWillDisappear} → \texttt{viewDidDisappear}; \texttt{deinit} last.
The \emph{app} moves Not running → Inactive \ensuremath{\rightleftarrows} Active → Background → Suspended, reported
per \textbf{scene} since iOS 13.}

\vspace{2pt}
\noindent\begin{tikzpicture}[sheet, x=1.49cm]
  \node[once] (n0) at (0,0) {init\\(coder: / nib:)};
  \node[once] (n1) at (1,0) {loadView};
  \node[once] (n2) at (2,0) {viewDid\\Load};
  \node[app]  (n3) at (3,0) {viewWill\\Appear};
  \node[app]  (n4) at (4,0) {viewIs\\Appearing};
  \node[lay]  (n5) at (5,0) {viewWill\\Layout\\Subviews};
  \node[lay]  (n6) at (6,0) {viewDid\\Layout\\Subviews};
  \node[app]  (n7) at (7,0) {viewDid\\Appear};
  \node[dis]  (n8) at (8,0) {viewWill\\Disappear};
  \node[dis]  (n9) at (9,0) {viewDid\\Disappear};
  \node[once, draw=sheetGrey, fill=black!6] (n10) at (10,0) {deinit};
  \foreach \a/\b in {n0/n1,n1/n2,n2/n3,n3/n4,n4/n5,n5/n6,n6/n7,n7/n8,n8/n9,n9/n10}
    \draw[flow] (\a) -- (\b);
  % braces above
  \draw[br, sheetBlue] (-0.45,0.55) -- node[lbl, above=2pt, text=sheetBlue]{\textbf{ONCE} per VC — first \texttt{.view} access} (2.45,0.55);
  \draw[br, sheetGreen!70!black] (2.55,0.55) -- node[lbl, above=2pt, pos=0.18, text=sheetGreen!60!black]{\textbf{EVERY} appearance} (7.45,0.55);
  \draw[br, sheetRed] (7.55,0.55) -- node[lbl, above=2pt, text=sheetRed]{\textbf{EVERY} disappearance} (9.45,0.55);
  \node[lbl, above=5pt, text=sheetGrey] at (10,0.55) {\textbf{ONCE}};
  % layout loop
  \draw[->, thick, sheetOrange] (n6.north) to[out=110,in=70] node[lbl, above, text=sheetOrange]{\textbf{MANY}} (n5.north);
  % what goes where
  \node[lbl, below=1pt] at (n0.south) {no \texttt{self.view}\\here};
  \node[lbl, below=1pt] at (n1.south) {\texttt{view = MyView()}\\no \texttt{super}};
  \node[lbl, below=1pt] at (n2.south) {subviews,\\constraints, bind};
  \node[lbl, below=1pt] at (n3.south) {refresh data,\\nav-bar style};
  \node[lbl, below=1pt] at (n4.south) {in hierarchy:\\traits + size OK};
  \node[lbl, below=1pt] at (5.5,-0.45) {each layout pass: rotation,\\bounds, safe area → frame maths};
  \node[lbl, below=1pt] at (n7.south) {start anim,\\video, timers};
  \node[lbl, below=1pt] at (n8.south) {save input,\\end editing};
  \node[lbl, below=1pt] at (n9.south) {stop timers,\\observers};
  \node[lbl, below=1pt] at (n10.south) {never runs\\if a cycle};
  % re-appear loop
  \draw[->, thick, sheetGreen!70!black, rounded corners=3pt]
    (n9.south) ++(0,-0.62) -- ++(0,-0.25) -- node[lbl, below, text=sheetGreen!60!black]
    {shown again (pop back to it, tab re-selected, \texttt{.fullScreen} modal dismissed) → \texttt{viewDidLoad} is NOT re-run}
    ($(n3.south)+(0,-0.87)$) -- ++(0,0.25);
\end{tikzpicture}

\begin{multicols}{2}
\raggedright

\section{How it works}
\begin{itemize}
  \item \textbf{\texttt{loadView}} — creates the root view. Default loads the nib/storyboard
        (or a blank \texttt{UIView}). Override \emph{only} for all-code UI; assign
        \texttt{view} and do \textbf{not} call \texttt{super}. Never call it yourself.
  \item \textbf{\texttt{viewDidLoad}} — in memory, \textbf{not} in a window, bounds = nib
        size: add subviews + constraints, no frame maths.
  \item \textbf{\texttt{viewIsAppearing(\_:)}} — iOS 17 SDK, \textbf{back-deployed to iOS 13}.
        After \texttt{viewWillAppear}: view is \emph{in the hierarchy}, trait collection and
        geometry are accurate. Best place for size/trait-dependent setup per appearance.
  \item \textbf{\texttt{viewWill/DidLayoutSubviews}} — every layout pass, also while on
        screen: keep cheap + idempotent. Always call \texttt{super}. Child VCs get callbacks
        forwarded (\texttt{addChild} → \texttt{addSubview} → \texttt{didMove(toParent:)}).
\end{itemize}

\section{Example}
\begin{lstlisting}[language=SwiftSheet]
final class FeedVC: UIViewController {
  let table = UITableView(); var timer: Timer?
  override func loadView() { view = table }  // no super
  override func viewDidLoad() { super.viewDidLoad(); setUp() }
  override func viewWillAppear(_ a: Bool) {
    super.viewWillAppear(a); reload() }      // every show
  override func viewDidLayoutSubviews() {    // many times
    super.viewDidLayoutSubviews(); gradient.frame = view.bounds }
  override func viewDidAppear(_ a: Bool) { super.viewDidAppear(a)
    timer = .scheduledTimer(withTimeInterval: 5, repeats: true) {
      [weak self] _ in self?.reload() } }
  override func viewWillDisappear(_ a: Bool) { // mirror
    super.viewWillDisappear(a); timer?.invalidate() }
  deinit { print("FeedVC freed") } }         // proves no leak
\end{lstlisting}

\section{App states — scene-based (iOS 13+)}
\begin{tikzpicture}[sheet]
  \node[st, draw=sheetGrey, fill=black!6] (nr) at (0,1.9) {Not running};
  \node[st] (in) at (3.0,1.9) {Inactive};
  \node[st, draw=sheetGreen!70!black, fill=sheetGreen!14] (ac) at (6.0,1.9) {Active};
  \node[st, draw=sheetOrange, fill=sheetOrange!10] (bg) at (3.0,0) {Background};
  \node[st, draw=sheetRed, fill=sheetRed!8] (su) at (0,0) {Suspended};
  \draw[flow] (nr) -- node[ed, above]{willConnectTo}node[ed,below]{launch} (in);
  \draw[flow, sheetGreen!70!black] (in.20) -- node[ed, above]{sceneDidBecome\\Active} (ac.160);
  \draw[flow, sheetOrange] (ac.200) -- node[ed, below]{sceneWillResign\\Active} (in.340);
  \draw[flow] (in.250) -- node[ed, left]{sceneDidEnter\\Background} (bg.110);
  \draw[flow] (bg.70) -- node[ed, right]{sceneWillEnter\\Foreground} (in.290);
  \draw[flow] (bg.190) -- node[ed, below]{no callback} (su.350);
  \draw[flow] (su.10) -- node[ed, above]{event wakes} (bg.170);
  \draw[hot, draw=sheetRed] (su) -- node[ed, right, text=sheetRed]{jetsam:\\silent kill} (nr);
  \node[lbl, anchor=west, text width=28mm, align=left] at (4.2,0.1) {Background: runs
     \textasciitilde30\,s (\texttt{beginBackgroundTask}) or a bg mode.
     Suspended: in RAM, \textbf{no code runs}.};
\end{tikzpicture}

\section{Launch sequence (cold)}
dyld + static init → \texttt{main} (\texttt{@main}) → \texttt{UIApplicationMain} →
\texttt{UIApplication} + delegate → \texttt{willFinishLaunching…} →
\texttt{didFinishLaunching…} → \texttt{configurationForConnecting} →
\texttt{scene(\_:willConnectTo:options:)} (make \texttt{UIWindow}, set root,
\texttt{makeKeyAndVisible}) → \texttt{sceneWillEnterForeground} →
\texttt{sceneDidBecomeActive} → first frame. Watchdog: \textasciitilde20\,s\unverified{} or killed
(\texttt{0x8badf00d}). A \textbf{warm} resume skips everything up to willEnterForeground.

\columnbreak

\section{App delegate vs scene delegate}
\textbf{\texttt{UIApplicationDelegate}} = the \emph{process}: \texttt{didFinishLaunching},
\texttt{configurationForConnecting}, \texttt{didDiscardSceneSessions}, push token,
\texttt{applicationWillTerminate} (unreliable).
\textbf{\texttt{UISceneDelegate}} = one \emph{UI instance} (N per app on iPad):
\texttt{willConnectTo}, the four hooks in the diagram, \texttt{sceneDidDisconnect}. With
scenes adopted, \texttt{application\{DidBecomeActive, WillResignActive, DidEnterBackground,
WillEnterForeground\}} are \textbf{not called}.
\textbf{SwiftUI:} \texttt{@Environment(\textbackslash.scenePhase)} = \texttt{.active} /
\texttt{.inactive} / \texttt{.background} (no ``suspended'' — no code runs); in
\texttt{App} it aggregates all scenes. Bridge with \texttt{@UIApplicationDelegateAdaptor}.

\section{Traps}
\begin{itemize}
  \trap{Touching \texttt{self.view} in \texttt{init} or a setter fires \texttt{loadView} +
        \texttt{viewDidLoad} early — before your properties/outlets are set.}
  \trap{Frames in \texttt{viewDidLoad} are the storyboard's, not the device's. Use
        \texttt{viewIsAppearing} / \texttt{viewDidLayoutSubviews}.}
  \trap{\texttt{.pageSheet} / \texttt{.automatic} modal (iOS 13 default): the
        \emph{presenter} gets \textbf{no} viewWillDisappear/viewWillAppear — refresh via a
        delegate or \texttt{presentationControllerDidDismiss}.}
  \trap{Pause a game on \texttt{WillResignActive} (control centre, call) — not on
        DidEnterBackground; save on DidEnterBackground. Never rely on willTerminate:
        jetsam kills a suspended app silently.}
  \trap{A strong \texttt{Timer}/closure target → \texttt{deinit} never runs. Stop it in
        \texttt{viewWillDisappear}.}
\end{itemize}

\section{Child view controllers — the exact calls}
\begin{lstlisting}[language=SwiftSheet]
addChild(child)                 // calls willMove(toParent:)
view.addSubview(child.view)     // + constraints
child.didMove(toParent: self)   // YOU call this one
// removing: the mirror image
child.willMove(toParent: nil)   // YOU call this one
child.view.removeFromSuperview()
child.removeFromParent()        // calls didMove(toParent: nil)
\end{lstlisting}
Appearance callbacks reach the child automatically
(\texttt{shouldAutomaticallyForwardAppearanceMethods} = \texttt{true}); a custom container
that turns it off calls \texttt{beginAppearanceTransition} / \texttt{endAppearanceTransition}.

\section{Remember}
\textbf{init · load · DidLoad · Will · Is · (layout)$^{n}$ · Did — Will · Did · deinit.}
\emph{Will} = last chance to prepare; \emph{Did} = already happened, safe to act.
Appear/disappear come in \textbf{pairs}: start in one, stop in its mirror.


\end{multicols}

\noindent{\footnotesize\color{sheetGrey}\textit{Related:} navigation \& presentation ·
ARC \& ownership (retain cycles) · Auto Layout (the layout pass) · background execution (\texttt{BGTaskScheduler}) ·
async work \& data flow in SwiftUI views (\texttt{onAppear}/\texttt{task}) · scenes, multiwindow + state restoration}

\end{document}
